% HD-S27. GPT-6 Astra Pro, /root/sixfolds, inherited effort not exposed.
% 21 November 2026. Exact scalar-Weyl scope theorem.
\subsection{HD-S27: every possible rank-nine natural \(E_6\) or \(E_7\) sixfold source is split}
\label{sec:e6-e7-split}

Consider connected sphere covers with full monodromy
\(G_m=W(E_m)\cdot\mu_4\), \(m=6,7\), and the natural
\(\mathbb Q(i)\)-representation. Their primitive Jacobian factors
and restricted theta forms are constructed by the explicit
matrix-coefficient idempotent below. Every such Weil sixfold with
a rank-nine period differential has ordinary Hermitian determinant
class \([-1]\). More precisely, all components with twelve branch
values have this determinant class, and all components with fewer
branch values have period rank at most eight. There are no components
with more than twelve branch values and primitive dimension six.
This is a statement about these full-group natural polarized factors.

\paragraph{1. Exact classes and the forced branch profiles.}
For \(m=6,7\), take the leading \(m\times m\) part \(C_m\)
of the Cartan matrix in HD-S22 and the root-basis simple reflections
\(1-u_ju_j^tC_m\). The independent finite computation
\texttt{weyl\_scalar\_class\_scan.g} starts with these integer
matrices, constructs their full root orbit, and computes all conjugacy
classes of their faithful root-permutation group. It imports no Weyl
class or character table. The complete class list, matrix for each
representative, class size and characteristic polynomial are saved
in \texttt{e6\_classes.json} and \texttt{e7\_classes.json}.
The results are
\[
\begin{array}{c|r|r|r|c|c}
 m&\#R&|W|&\#\text{classes}&-I\in W&
       \max_w\dim\ker(w-i)\\\hline
 6&72&51840&25&\text{no}&2\\
 7&126&2903040&60&\text{yes}&2
\end{array}
\]
For \(E_6\), the unique class attaining the maximum has polynomial
\(\Phi_4^2(X-1)^2\). For \(E_7\), the four such classes have
\[
 \Phi_4^2(X-1)^{3-r}(X+1)^r,\qquad r=0,1,2,3.
\]
All these representatives have order four. The scan also checks that
the only codimension-one elements of \(G_6\) are the Weyl
reflections: no \(-w\), \(w\in W(E_6)\), has fixed dimension
five, and neither scalar-odd coset can contain such an element.
For \(G_7\) the real cosets coincide, and the only codimension-one
class is again the ordinary Weyl-reflection class. The representation
bridge is explicit: if \(p(M)\) is the permutation of column roots,
then \(p(MN)=p(N)p(M)\), and the genuine representation is
\(p(M)\mapsto M^t\). It is conjugate to the original root-basis
representation by \(C_m\).

The quotient \(G_6/W(E_6)\) is cyclic of order four, and
\(G_7/W(E_7)\) is cyclic of order two. A generating product-one
tuple has at least two scalar-odd entries. These have codimension
at least four for \(E_6\), and at least five for \(E_7\).
The exact parabolic quotient of HD-S23 gives
\(\sum c_j=2m+6\) for a sixfold. Thus in both cases
\[
 b\le12.
\]
Equality forces exactly two minimal scalar-odd entries and ten Weyl
reflections. In the \(E_6\) case the scalar phases of the two
exceptional entries are \(i\) and \(-i\), because their product
in the order-four scalar quotient is one. In every equality case
all inertia orders divide four. As in HD-S26, the period rank is at
most \(b-3\), with the exact marked-Hurwitz and Jacobian maps.
The same codimension argument gives \(b-3\le11\) for natural
Weil eightfolds from either group.

\paragraph{2. Integral parabolic duality with every map retained.}
Put \(R=\mathbb Z[i,1/2]\), a principal ideal domain. The following
argument applies to each equality-case tuple above. Let \(\Sigma\)
be the sphere with disjoint open branch discs removed. Let \(L\)
be a free \(R\)-coefficient lattice with the natural monodromy,
and \(L^\vee=\operatorname{Hom}_R(L,R)\) its dual local system.
The central scalar \(iI\) belongs to the generated group, and
\(i-1\) is a unit in \(R\). Since any word \(g\) satisfies
\((g-1)L\subseteq\sum_j(g_j-1)L\), this proves
\[
 H_0(\Sigma,L)=H_0(\Sigma,L^\vee)=0.
\]
The ordinary first homology is free: the surface retracts onto a
finite graph, and it is the kernel in the resulting complex of free
modules. Poincar\'e--Lefschetz duality and the universal coefficient
sequence give perfect cross-pairings
\[
 H_1(\Sigma,L^\vee)\ \times\
 H_1(\Sigma,\partial\Sigma;L)\longrightarrow R,
\]
and the version with the two coefficient systems interchanged. In
particular the relative first homology is free as well. The exact
duality and coefficient maps are the cap-product maps; a primary
source over a commutative coefficient ring is Dettweiler--Wewers,
Proposition 2.1, pp.153--154, together with their universal-coefficient
discussion on p.153.

For a boundary circle with inertia \(g\) of order dividing four,
the averaging idempotent \(\frac1{\operatorname{ord}g}\sum g^j\)
identifies its coinvariants and invariants as a direct summand of
\(L\). Therefore both boundary homology modules are free. The
long exact sequence of the pair shows that the image
\[
 P(L)=\operatorname{im}\bigl(H_1(\Sigma,L)
                 \longrightarrow H_1(\Sigma,\partial\Sigma;L)\bigr)
\]
has free cokernel \(H_0(\partial\Sigma,L)\); it is a primitive
submodule of relative homology. Naturality of the cap product makes
the two absolute-to-relative maps adjoint, with the usual alternating
sign. It follows that the induced pairing
\[
                       P(L^\vee)\times P(L)\longrightarrow R
\]
is perfect, not merely nondegenerate over the fraction field. Indeed
a functional on the primitive submodule \(P(L)\) extends to the
whole relative module. Its unique dual absolute class gives the
required element of \(P(L^\vee)\); the kernel is exactly the
kernel of the absolute-to-relative map, by adjointness and the
perfect cross-pairing. This proves both injectivity and surjectivity.

Equivalently one may use \(H^1(\mathbb P^1,j_*L)\), the integral
parabolic cohomology. Here its identification with the image just
used is the restriction map followed by Poincar\'e duality. To check
this with integral coefficients, cover the sphere by \(\Sigma\)
and the branch discs. On each disc the extension is the invariant
summand of a finite branched-disc pushforward, with averaging degree
dividing four. It has no higher cohomology. The Mayer--Vietoris
sequence then identifies its first cohomology with the kernel of
restriction to the boundary circles, which is the image of compact
support cohomology. Thus these identifications use actual maps and
preserve the cup pairing.

\paragraph{3. The \(E_7\) determinant.}
The natural root lattice has Gram matrix \(C_7\), with
\(\det C_7=2\). It is therefore Hermitian self-dual over \(R\).
The perfect parabolic pairing in paragraph 2 becomes a perfect
skew-Hermitian form on the rank-six parabolic lattice. Multiplication
by \(i\) gives the Hermitian cup form used in HD-S23, still perfect.
Its determinant belongs to \(R^\times\cap\mathbb Q\), and hence
is \(\pm2^a\) for an integer \(a\). Its signature \((3,3)\)
makes it negative. Since \(2=N_{K/\mathbb Q}(1+i)\), its ordinary
determinant class is exactly \([-1]\).

\paragraph{4. The \(E_6\) discriminant module disappears by an exact map.}
The root lattice now has Gram matrix \(C_6\), of determinant 3.
Write \(L^\#\) for its Hermitian dual inside the same rational
representation. Then
\[
 0\longrightarrow L\longrightarrow L^\#
   \longrightarrow A\longrightarrow0,
 \qquad A\simeq R/(3).
\]
Every Weyl reflection acts trivially on \(A\): for \(x\in L^\#\)
and a norm-two root \(\alpha\),
\(s_\alpha x-x=-h(\alpha,x)\alpha\in L\). The scalar \(i\)
acts as multiplication by \(i\). Consequently the discriminant
local system of every equality-case tuple has precisely two
nontrivial branch monodromies, \(i\) and \(-i\); the ten
reflection punctures have trivial monodromy on \(A\).

After filling those ten trivial punctures the remaining local system
is on \(\mathbb C^*\), with its cellular complex given by
\(i-1:A\to A\). This map is an isomorphism, so its ordinary
cohomology vanishes. The compact-support cohomology also vanishes by
duality over the field \(R/(3)\simeq\mathbb F_9\). At the two
remaining branch points the invariant stalk is zero. Therefore
\[
 H^0(\mathbb P^1,j_*A)=H^1(\mathbb P^1,j_*A)=0.
\]
Taking local invariants in the displayed coefficient exact sequence
is exact: all inertia groups have order dividing four, a unit in
\(R\). Hence it gives an exact sequence of \(j_*\)-sheaves, and
the cohomology sequence gives the actual isomorphism
\[
 H^1(\mathbb P^1,j_*L)
       \xrightarrow{\ \sim\ }
 H^1(\mathbb P^1,j_*L^\#).
\]
This is induced by inclusion of the coefficient lattices. Combining
it with the perfect cross-duality of paragraph 2 makes the parabolic
Hermitian cup lattice self-dual over \(R\). Its determinant is
again a negative rational unit of \(R\), hence has class \([-1]\).
The assertion uses the exact discriminant local system and its
vanishing, rather than treating the coefficient determinant 3 as a
harmless change of scale.

\paragraph{5. The actual primitive polarized factor.}
For completeness, the map from these coefficient forms to the
Jacobian is precisely the construction of HD-S22--HD-S23 with 8
replaced by \(m\). Take a norm-two root \(v\), and put
\[
 e=\frac m{|G_m|}\sum_{g\in G_m}
       \operatorname{Re}h_0(v,\rho(g^{-1})v)\,g,
 \qquad D_m=|G_m|/m.
\]
The connected root diagram gives scalar commutant, so the same
matrix-coefficient averaging proves \(e^2=e=e^*\) and
\(e\mathbb Q[G_m]e=K\). Root integrality gives
\(D_me\in\mathbb Z[G_m]\). Here
\[
 |G_6|=207360,\quad D_6=34560,\qquad
 |G_7|=5806080,\quad D_7=829440.
\]
Use the covariant Jacobian action \(a_g=L_{g,*}\), let
\(B=\operatorname{im}(D_me)\), and choose the integral Gaussian
action \(i_B=a_{iI}^{-1}\). The coefficient function
\(j_x(g)=h_0(v,\rho(g)x)/2\) satisfies
\[
 \sum_g\bar j_x(g)j_y(g)=\frac{|G_m|}{2m}h_0(x,y),
 \qquad h_B(\Phi x,\Phi y)=D_m H(x,y).
\]
These are the same actual Poincar\'e-duality and Abel maps as in
HD-S23, with the action convention stated there. Multiplication of
a rank-six determinant by \(D_m^6=N(D_m^3)\) preserves its norm
class. The change from root coordinates to the fundamental-weight
coordinates is the explicit map \(x\mapsto C_mx\); it intertwines
the simple-reflection matrices and pulls \(C_m^{-1}\) back to
\(C_m\). It too changes the Gram determinant by a norm.
This proves the stated determinant conclusion for the actual
primitive polarized Jacobian factors.

\paragraph{Checks and scope.}
The whole proof above was read by its author. The exact class scan
is exhaustive within the two explicitly generated finite groups.
Eight separately constructed generating twelve-point tuples for
each group are recorded in \texttt{e6\_sixfold\_candidates.json}
and \texttt{e7\_sixfold\_candidates.json}, with full parabolic cup
matrices; all have class \([-1]\). These examples check the formula
and nonempty branch profiles, but the universal conclusion follows
from paragraphs 1--5, not from sampling. No period-rank-nine equality
or new nonsplit coverage is asserted for those examples.
